\documentclass[12pt,a4paper]{report}
\usepackage[utf8]{inputenc}
\usepackage[T1]{fontenc}
\usepackage[french]{babel}
\usepackage{geometry}
\geometry{left=3cm,right=2.5cm,top=3cm,bottom=3cm}
\usepackage{graphicx}
\usepackage{booktabs}
\usepackage{float}
\usepackage{hyperref}
\usepackage{listings}
\usepackage{xcolor}
\usepackage{amsmath}
\usepackage{setspace}
\usepackage{cite}
% CORRECTION PROF : legendes homogenes
\usepackage{caption}
\captionsetup{
  font=small,
  labelfont=bf,
  textfont=it,
  justification=centering
}
% CORRECTION PROF : charte graphique en-tete/pied de page
\usepackage{fancyhdr}
\pagestyle{fancy}
\fancyhf{}
\fancyhead[L]{\small\textit{Architecture SDN -- UASZ 2024-2025}}
\fancyhead[R]{\small\textit{Aminata Thior}}
\fancyfoot[C]{\thepage}
\renewcommand{\headrulewidth}{0.4pt}
\onehalfspacing

\lstset{
  basicstyle=\ttfamily\small,
  keywordstyle=\color{blue}\bfseries,
  commentstyle=\color{green!60!black},
  stringstyle=\color{orange},
  numbers=left,
  numberstyle=\tiny\color{gray},
  frame=single,
  breaklines=true,
  captionpos=b,
  showstringspaces=false,
  tabsize=4
}

\begin{document}

\begin{titlepage}
\begin{center}
{\large\bfseries REPUBLIQUE DU SENEGAL}\\[0.3cm]
\includegraphics[width=0.25\textwidth]{drapeau.png}\\[0.3cm]
{\large\textcolor{green}{\textbf{Un peuple}} \_
\textcolor{yellow}{\textbf{Un but}} \_
\textcolor{red}{\textbf{Une foi}}}\\[0.3cm]
{\large **************}\\[0.3cm]
\rule{\linewidth}{0.3mm}\\[0.2cm]
{\small\bfseries Ministère de l'Enseignement Supérieur,\\
de la Recherche et de l'Innovation}\\[0.2cm]
\rule{\linewidth}{0.3mm}\\[0.2cm]
{\large\bfseries Université Assane Seck de Ziguinchor}\\[0.2cm]
{\large Faculté des Sciences et Technologies}\\[0.2cm]
{\large Département d'Informatique}\\[0.5cm]
\rule{\linewidth}{0.5mm}\\[0.3cm]
{\large Mémoire de Master 1 -- Réseaux Avancés}\\[0.5cm]
{\large\bfseries Thème :}\\[0.3cm]
{\huge\bfseries Architecture SDN pour la Gestion\\[0.3cm]
Dynamique du Trafic et la\\[0.3cm]
Sécurité Réseau}\\[0.5cm]
\rule{\linewidth}{0.5mm}\\[0.5cm]
% CORRECTION PROF : alignement page de garde avec tableau
\begin{tabular}{p{0.45\textwidth}p{0.45\textwidth}}
\textbf{Auteure :} & \textbf{Encadrant :} \\
Aminata Thior & Pr. Malick Ndoye \\
\end{tabular}\\[1cm]
{\large Année universitaire 2024--2025}
\end{center}
\end{titlepage}

% Page de titre sans en-tete ni pied de page
\thispagestyle{empty}

\chapter*{Remerciements}
\addcontentsline{toc}{chapter}{Remerciements}
Je tiens à exprimer ma profonde gratitude à toutes les
personnes qui ont contribué à la réalisation de ce projet.

En premier lieu, je remercie sincèrement mon encadrant,
\textbf{Professeur Malick Ndoye}, pour sa disponibilité,
ses conseils avisés et son soutien tout au long de ce
travail. Sa rigueur scientifique et ses orientations ont
été déterminantes dans la réussite de ce projet.

Je remercie également l'ensemble du corps enseignant du
\textbf{Département d'Informatique} de l'Université Assane
Seck de Ziguinchor pour la qualité de la formation dispensée
tout au long de mon cursus de Master.

Je tiens à exprimer ma reconnaissance à ma \textbf{famille},
pour leur soutien indéfectible, leur patience et leurs
encouragements constants durant toutes mes années d'études.

Enfin, je remercie mes camarades de promotion pour les
échanges enrichissants et l'entraide tout au long de
cette année universitaire.
\newpage

\chapter*{Résumé}
\addcontentsline{toc}{chapter}{Résumé}
Les réseaux traditionnels souffrent de rigidité et d'un
manque de visibilité globale, rendant difficile la gestion
dynamique du trafic et la réponse rapide aux menaces de
sécurité. Le paradigme SDN (Software-Defined Networking)
propose une approche innovante en séparant le plan de
contrôle du plan de données, offrant ainsi une vue
centralisée et programmable du réseau.

Ce projet présente la conception et le déploiement d'une
architecture SDN complète basée sur le contrôleur Ryu et
l'émulateur Mininet. Quatre modules ont été développés :
un module de surveillance du trafic en temps réel, un module
de routage dynamique basé sur les métriques de performance,
un module QoS pour la priorisation des flux, et un module
de sécurité intégrant la détection DDoS par entropie de
Shannon, la mitigation automatique graduée et un pare-feu
stateful.

Les évaluations sur une topologie arbre (16 hôtes) montrent
un débit TCP de 23,5 Gbps, une latence de 0,134 ms, une
détection DDoS en 5 secondes et une réduction de la latence
de 52\% par rapport à un réseau traditionnel.

\textbf{Mots-clés :} SDN, OpenFlow, Ryu, Mininet, DDoS,
QoS, Sécurité réseau, Gestion du trafic.
\newpage

\chapter*{Abstract}
\addcontentsline{toc}{chapter}{Abstract}
Traditional networks suffer from rigidity and lack of global
visibility, making dynamic traffic management and rapid
response to security threats difficult. The SDN paradigm
proposes an innovative approach by separating the control
plane from the data plane, offering a centralized and
programmable view of the network.

This project presents the design and deployment of a complete
SDN architecture based on the Ryu controller and Mininet.
Four modules were developed: real-time traffic monitoring,
dynamic routing, QoS flow prioritization, and a security
module with DDoS detection via Shannon entropy, automatic
mitigation and a stateful firewall.

Evaluations show 23.5 Gbps TCP throughput, 0.134 ms latency,
DDoS detection in 5 seconds and 52\% latency reduction.

\textbf{Keywords:} SDN, OpenFlow, Ryu, Mininet, DDoS, QoS,
Network Security, Traffic Management.
\newpage

\tableofcontents
\newpage
\listoffigures
\addcontentsline{toc}{chapter}{Liste des figures}
\newpage
\listoftables
\addcontentsline{toc}{chapter}{Liste des tableaux}
\newpage

\chapter*{Glossaire}
\addcontentsline{toc}{chapter}{Glossaire}

\textbf{ACL (Access Control List)} : Liste de règles de
filtrage définissant les flux autorisés ou refusés. Dans
ce projet, les ACL bloquent Telnet (23), FTP (21), RPC
(135), TFTP (69) et autorisent SSH (22), HTTP (80),
HTTPS (443), iperf3 (5201) et ICMP.\\[0.3cm]

\textbf{DDoS (Distributed Denial of Service)} : Attaque
par déni de service distribué visant à saturer un réseau
en envoyant un volume massif de trafic depuis plusieurs
sources simultanées. Les types courants comprennent le
SYN flood (TCP), l'UDP flood et l'ICMP flood. Dans ce
projet, ce type d'attaque est détecté en 5 secondes.\\[0.3cm]

\textbf{Entropie de Shannon} : Mesure statistique de la
diversité du trafic réseau définie par la formule
$H = -\sum p_i \log_2(p_i)$, normalisée entre 0 et 1.
Une entropie proche de 0 indique un trafic concentré
(attaque probable), tandis qu'une valeur proche de 1
correspond à un trafic normal et uniformément distribué.
Le seuil de détection utilisé dans ce projet est 0.5.\\[0.3cm]

\textbf{FlowMod} : Message OpenFlow (OFPFlowMod) envoyé
par le contrôleur aux switches pour installer, modifier
ou supprimer des règles de flux dans la table de flux.
Ses paramètres clés sont : match, priority, actions,
idle\_timeout et hard\_timeout.\\[0.3cm]

\textbf{Honeypot (Pot de miel)} : Système leurre conçu
pour attirer les attaquants afin d'analyser leur
comportement. Dans ce projet, h16 (10.0.0.16) joue le
rôle de honeypot. Le trafic suspect y est redirigé via
l'action SetField d'OpenFlow (priorité 90, timeout 120s).\\[0.3cm]

\textbf{Mininet} : Émulateur de réseau open-source
permettant de créer des topologies réseau virtuelles
complètes sur une seule machine Linux. Il utilise les
namespaces Linux et Open vSwitch pour émuler switches,
hôtes et liens. La version 2.3+ est utilisée dans
ce projet.\\[0.3cm]

\textbf{Mitigation} : Ensemble de mesures automatiques
prises pour réduire l'impact d'une attaque réseau. Ce
projet implémente trois niveaux progressifs : Rate
Limiting (priorité 80, 30s), Honeypot (priorité 90,
120s) et DROP total (priorité 100, 60s).\\[0.3cm]

\textbf{OpenFlow} : Protocole de communication standard
entre le contrôleur SDN et les switches réseau. Il définit
une table de flux composée de règles spécifiant les
critères de correspondance (match), la priorité, les
actions, les compteurs et les timeouts. La version 1.3,
utilisée dans ce projet, introduit les tables multiples
et les statistiques améliorées.\\[0.3cm]

\textbf{OVS (Open vSwitch)} : Switch virtuel open-source
supportant OpenFlow. Il implémente un datapath kernel
Linux hautes performances et est utilisé par Mininet
pour émuler les switches OpenFlow dans ce projet.\\[0.3cm]

\textbf{PacketIn} : Message OpenFlow envoyé par un switch
au contrôleur lorsqu'un paquet ne correspond à aucune
règle de flux (table-miss). Ce message déclenche le
traitement par le contrôleur, qui installe ensuite une
nouvelle règle via FlowMod.\\[0.3cm]

\textbf{QoS (Quality of Service)} : Ensemble de mécanismes
garantissant des niveaux de performance (débit, latence,
jitter) pour certains types de flux réseau. Ce projet
implémente trois niveaux : HAUTE (iperf3 TCP 5201, SSH),
MOYENNE (TCP général, ICMP) et BASSE (UDP).\\[0.3cm]

\textbf{Ryu} : Framework de contrôleur SDN open-source
développé en Python par NTT Labs. Il supporte OpenFlow
1.0 à 1.5 et fournit une API événementielle simple basée
sur des décorateurs Python. Il est utilisé dans ce projet
pour l'implémentation de tous les modules.\\[0.3cm]

\textbf{SDN (Software-Defined Networking)} : Paradigme
réseau séparant le plan de contrôle du plan de données.
Un contrôleur centralisé dispose d'une vue globale du
réseau et programme dynamiquement les switches via
OpenFlow. Ce paradigme a été introduit en 2008 à
l'Université de Stanford par McKeown et al.\\[0.3cm]

\textbf{Topologie Fat-Tree} : Architecture réseau
datacenter organisée en trois niveaux (Core, Agrégation,
Edge) offrant une bande passante bisectionnelle maximale
et une haute disponibilité grâce aux chemins alternatifs.
Pour k=4, elle comprend 16 hôtes, 4 pods, 4 switches
core, 8 switches agrégation et 8 switches edge.
\newpage

\chapter{Introduction Générale}

\section{Contexte et Motivation}

Les réseaux informatiques constituent aujourd'hui l'épine
dorsale de toute infrastructure numérique mondiale. En
Afrique de l'Ouest et particulièrement au Sénégal, le
développement des infrastructures numériques représente
un enjeu stratégique majeur dans le cadre du programme
``Sénégal Numérique 2025''. Avec la multiplication des
universités connectées et l'essor des services numériques,
la gestion efficace et sécurisée des réseaux informatiques
est devenue une priorité nationale \cite{ndiaye2022sdn}.

Dans ce contexte, les architectures réseau traditionnelles
montrent leurs limites face aux exigences croissantes de
flexibilité, d'agilité et de sécurité. En effet, ces
réseaux présentent plusieurs limitations majeures qu'il
convient d'analyser :

\begin{itemize}
    \item \textbf{Rigidité} : La configuration manuelle
    de chaque équipement rend le réseau difficile à
    adapter dynamiquement aux variations de trafic, ce
    qui entraîne des délais importants lors de changements
    de politique réseau.
    \item \textbf{Complexité de gestion} : L'administration
    de réseaux hétérogènes nécessite une expertise
    spécifique pour chaque type d'équipement, sans vue
    globale cohérente de l'ensemble du réseau.
    \item \textbf{Manque de visibilité globale} : L'absence
    d'une vue centralisée complique la détection d'anomalies
    et l'optimisation des performances, car chaque
    équipement n'a qu'une vision locale du trafic.
    \item \textbf{Sécurité réactive} : Les politiques de
    sécurité sont configurées statiquement et ne s'adaptent
    pas en temps réel aux menaces émergentes, laissant le
    réseau vulnérable pendant les délais de réaction.
\end{itemize}

Le paradigme SDN (Software-Defined Networking) propose de
séparer le plan de contrôle du plan de données via un
contrôleur centralisé disposant d'une vision globale du
réseau et pouvant programmer dynamiquement les règles de
transfert sur les commutateurs via le protocole OpenFlow
\cite{mckeown2008openflow}.

\section{Problématique}

\textbf{Comment concevoir une architecture SDN capable de
gérer dynamiquement le trafic réseau tout en intégrant des
mécanismes de détection et de mitigation automatique des
attaques réseau ?}

Cette problématique soulève plusieurs questions spécifiques
auxquelles ce projet apporte des réponses :
\begin{enumerate}
    \item Comment surveiller en temps réel les flux réseau
    dans un environnement SDN ?
    \item Comment adapter dynamiquement le routage en
    fonction des métriques de performance ?
    \item Comment détecter et atténuer automatiquement les
    attaques DDoS par analyse statistique du trafic ?
    \item Quelles sont les performances d'une architecture
    SDN par rapport à un réseau traditionnel dans un
    contexte académique ?
\end{enumerate}

\section{Objectifs du Projet}
\begin{description}
    \item[O1] Concevoir une architecture SDN complète
    identifiant les composants clés et leurs interactions.
    \item[O2] Déployer un environnement de simulation
    reproductible avec Mininet et le contrôleur Ryu.
    \item[O3] Implémenter un module de gestion dynamique
    du trafic comprenant le monitoring, le routage adaptatif,
    la QoS et le load balancing.
    \item[O4] Développer un module de sécurité réseau
    intégrant la détection DDoS, la mitigation graduée et
    un pare-feu stateful.
    \item[O5] Évaluer les performances à travers des
    scénarios de test réalistes et comparer avec un réseau
    traditionnel simulé dans le même environnement.
\end{description}

\section{Structure du Rapport}
Ce rapport est organisé en six chapitres complémentaires :
\begin{itemize}
    \item \textbf{Chapitre 2} : Présente l'état de l'art
    sur le SDN, OpenFlow, les contrôleurs et les mécanismes
    de sécurité réseau.
    \item \textbf{Chapitre 3} : Décrit l'architecture
    proposée et les quatre modules développés.
    \item \textbf{Chapitre 4} : Détaille l'implémentation
    de chaque module avec des extraits de code Python.
    \item \textbf{Chapitre 5} : Présente les tests,
    l'évaluation des performances et la comparaison avec
    un réseau traditionnel.
    \item \textbf{Chapitre 6} : Conclut le rapport et
    propose des perspectives d'amélioration.
\end{itemize}
\newpage

\chapter{État de l'Art}

\section{Introduction}
Ce chapitre présente une revue de la littérature sur le
Software-Defined Networking (SDN), les protocoles associés,
les contrôleurs existants, les mécanismes de sécurité et
la gestion de la qualité de service dans les réseaux
programmables. Cette étude bibliographique a guidé les
choix architecturaux et technologiques de ce projet.

\section{Le Paradigme SDN}

\subsection{Historique et Origine}
Le concept de Software-Defined Networking a émergé à
l'Université de Stanford autour de 2008, avec les travaux
fondateurs de McKeown et al. \cite{mckeown2008openflow}
qui ont introduit le protocole OpenFlow. Avant l'avènement
du SDN, les réseaux étaient gérés de manière entièrement
distribuée : chaque équipement intégrait son propre plan
de contrôle, rendant la gestion complexe et peu flexible
face aux besoins croissants des infrastructures modernes.

L'idée centrale du SDN consiste à séparer le plan de
contrôle du plan de données, ce qui permet une
programmation centralisée et flexible du réseau depuis
un contrôleur unique disposant d'une vue globale.
Selon Kreutz et al. \cite{kreutz2015software}, le SDN
repose sur quatre piliers fondamentaux :
\begin{enumerate}
    \item Les plans de données et de contrôle sont
    découplés et indépendants.
    \item Les décisions de transfert des paquets sont
    basées sur les flux et non sur les destinations.
    \item La logique de contrôle est externalisée vers
    un contrôleur accessible via une API ouverte.
    \item Le réseau est entièrement programmable par des
    applications externes au-dessus du contrôleur.
\end{enumerate}

\subsection{Architecture SDN en Couches}
L'architecture SDN se décompose en trois couches distinctes
interconnectées par des interfaces standardisées :

\begin{enumerate}
    \item \textbf{Couche Infrastructure (Data Plane)} :
    Cette couche regroupe les équipements réseau physiques
    ou virtuels (switches, routeurs) qui transfèrent les
    paquets selon les règles reçues du contrôleur. Dans
    notre projet, elle est émulée par Open vSwitch via
    Mininet.
    \item \textbf{Couche Contrôle (Control Plane)} :
    Le contrôleur SDN constitue le cerveau du réseau.
    Il maintient une vue globale et cohérente de la
    topologie, calcule les chemins optimaux et programme
    dynamiquement les règles de flux via l'interface
    southbound (OpenFlow). Le contrôleur Ryu joue ce
    rôle dans notre projet.
    \item \textbf{Couche Application (Application Plane)} :
    Les applications réseau (QoS, sécurité, monitoring,
    load balancing) interagissent avec le contrôleur via
    l'interface northbound (REST API, Python API). Les
    modules Python développés dans ce projet constituent
    cette couche applicative.
\end{enumerate}

\begin{table}[H]
\centering
\caption{Architecture en couches SDN et composants utilisés}
\label{tab:couches_sdn}
\begin{tabular}{llll}
\toprule
\textbf{Couche} & \textbf{Rôle} &
\textbf{Interface} & \textbf{Notre projet} \\
\midrule
Application & Logique métier & Northbound API &
Modules Python \\
Contrôle & Gestion centralisée & OpenFlow (SB) &
Ryu 4.34 \\
Infrastructure & Transfert paquets & OpenFlow (SB) &
OVS + Mininet \\
\bottomrule
\end{tabular}
\end{table}

\subsection{Avantages du SDN par rapport aux Réseaux Traditionnels}

\begin{table}[H]
\centering
\caption{Comparaison SDN vs réseau traditionnel}
\label{tab:sdn_vs_trad}
\begin{tabular}{lll}
\toprule
\textbf{Critère} & \textbf{Réseau Traditionnel} & \textbf{SDN} \\
\midrule
Contrôle & Distribué par équipement & Centralisé \\
Configuration & Manuelle (CLI) & Programmable (API) \\
Visibilité & Locale & Globale \\
Adaptabilité & Statique & Dynamique \\
Sécurité & Réactive & Proactive \\
Déploiement & Lent & Rapide \\
\bottomrule
\end{tabular}
\end{table}

\section{Le Protocole OpenFlow}

\subsection{Historique et Principe}
OpenFlow a été proposé par McKeown et al. en 2008 comme
premier protocole southbound standard pour le SDN
\cite{openflow_spec}. Il permet au contrôleur de programmer
les tables de flux des switches de manière uniforme,
indépendamment du constructeur de l'équipement.

\subsection{Structure d'une Table de Flux}
Une table de flux OpenFlow est composée d'entrées
structurées selon les champs suivants :

\begin{table}[H]
\centering
\caption{Structure d'une entrée de flux OpenFlow}
\label{tab:flow_entry}
\begin{tabular}{ll}
\toprule
\textbf{Champ} & \textbf{Description} \\
\midrule
Match Fields & Critères d'identification du paquet \\
Priority & Ordre de traitement (plus élevé = prioritaire) \\
Counters & Compteurs de paquets et d'octets \\
Instructions & Actions à appliquer au paquet \\
Timeouts & idle\_timeout et hard\_timeout \\
Cookie & Identifiant de la règle côté contrôleur \\
\bottomrule
\end{tabular}
\end{table}

\subsection{Champs de Correspondance et Actions}
OpenFlow 1.3 supporte de nombreux champs de correspondance
aux couches 2 (MAC, type Ethernet), 3 (IP src/dst,
protocole) et 4 (ports TCP/UDP). Les principales actions
disponibles sont : Output (envoi vers un port),
Drop (suppression du paquet par liste d'actions vide),
SetField (modification d'un champ comme la réécriture IP
pour le honeypot) et Group (groupe d'actions multipath).

\subsection{Évolution des Versions OpenFlow}

\begin{table}[H]
\centering
\caption{Évolution des versions OpenFlow}
\label{tab:of_versions}
\begin{tabular}{lll}
\toprule
\textbf{Version} & \textbf{Année} & \textbf{Nouveautés} \\
\midrule
1.0 & 2009 & Table unique, actions basiques \\
1.1 & 2011 & Tables multiples, groupes \\
1.2 & 2011 & Match étendu, IPv6 \\
1.3 & 2012 & Meters, tunnels, améliorations stats \\
1.4 & 2013 & Synchronisation, bundles \\
1.5 & 2014 & Egress tables, statistiques avancées \\
\bottomrule
\end{tabular}
\end{table}

Ce projet utilise \textbf{OpenFlow 1.3} pour sa stabilité,
son support des statistiques avancées (port\_stats,
flow\_stats) et sa large adoption dans la communauté
académique \cite{openflow_spec}.

\subsection{Fonctionnement du Pipeline OpenFlow}
Lorsqu'un paquet arrive sur un switch OpenFlow, celui-ci
consulte sa table de flux. Si une règle correspond au
paquet (match), les actions associées sont immédiatement
appliquées. Si aucune règle ne correspond (table-miss),
le switch envoie le paquet au contrôleur via un message
PacketIn. Le contrôleur analyse alors le paquet, prend
une décision de routage ou de sécurité, et installe une
nouvelle règle de flux via FlowMod. Les paquets suivants
appartenant au même flux sont ensuite traités directement
par le switch sans solliciter le contrôleur.

\section{Contrôleurs SDN Open-Source}

\subsection{Présentation des Principaux Contrôleurs}

\subsubsection{Ryu}
Ryu est un framework de contrôleur SDN open-source
développé en Python par NTT Labs \cite{ryu_doc}. Il
supporte OpenFlow 1.0 à 1.5 et fournit une API
événementielle simple basée sur des décorateurs Python,
ce qui facilite le développement rapide d'applications
réseau personnalisées.

\subsubsection{ONOS}
ONOS (Open Network Operating System) est un contrôleur
SDN haute performance développé en Java par l'Open
Networking Foundation (ONF). Il est conçu pour les
déploiements en production avec haute disponibilité,
scalabilité et support multi-contrôleurs.

\subsubsection{OpenDaylight}
OpenDaylight est un contrôleur SDN open-source soutenu
par la Linux Foundation. Son architecture modulaire via
OSGi/Karaf supporte de nombreux protocoles southbound
(OpenFlow, NETCONF, OVSDB), mais sa complexité de
configuration en limite l'usage académique.

\subsubsection{Floodlight}
Floodlight est un contrôleur SDN Java développé par Big
Switch Networks. Il propose une API REST simple et un
bon support d'OpenFlow 1.0 à 1.3, mais son développement
a ralenti ces dernières années.

\subsection{Tableau Comparatif}

\begin{table}[H]
\centering
\caption{Comparaison détaillée des contrôleurs SDN}
\label{tab:controleurs_detail}
\begin{tabular}{lllll}
\toprule
\textbf{Critère} & \textbf{Ryu} & \textbf{ONOS} &
\textbf{ODL} & \textbf{Floodlight} \\
\midrule
Langage & Python & Java & Java & Java \\
OpenFlow & 1.0-1.5 & 1.0-1.5 & 1.0-1.5 & 1.0-1.3 \\
Scalabilité & Moyenne & Haute & Haute & Moyenne \\
Complexité & Faible & Haute & Très haute & Moyenne \\
Communauté & Académique & Industrielle & Industrielle & Mixte \\
Intégration Mininet & Excellente & Bonne & Bonne & Bonne \\
\bottomrule
\end{tabular}
\end{table}

\subsection{Justification du Choix de Ryu}
Ryu a été retenu pour ce projet pour plusieurs raisons
complémentaires. Sa simplicité d'implémentation, grâce à
une API Python basée sur des décorateurs, permet un
développement rapide et lisible. Son support complet
d'OpenFlow 1.0 à 1.5, son intégration native avec Mininet
et sa large adoption dans la communauté académique en font
un choix idéal pour le prototypage. Enfin, son faible
overhead (21,5\% CPU, 65 MB mémoire) le rend adapté aux
environnements contraints \cite{benzekki2016software}.

\section{Topologies Réseau SDN}

\subsection{Topologie en Arbre}
La topologie en arbre est la plus utilisée dans les
environnements SDN académiques. Elle se caractérise par
un switch racine connecté à des switches intermédiaires,
auxquels sont reliés les hôtes finaux. Si elle présente
l'avantage d'une structure simple et facile à administrer,
elle souffre d'un inconvénient majeur : il n'existe qu'un
seul chemin entre toute paire source-destination, ce qui
la rend vulnérable aux pannes de liens.

Dans ce projet, la topologie arbre depth=2, fanout=4
comprend 5 switches et 16 hôtes, offrant un bon compromis
entre complexité et représentativité.

\subsection{Topologie Fat-Tree}
La topologie Fat-Tree, proposée par Al-Fares et al.
\cite{alfares2008scalable}, constitue une architecture
datacenter à trois niveaux (Core, Agrégation, Edge) qui
offre une bande passante bisectionnelle maximale, des
chemins multiples $(k/2)^2$ entre tout pod source et
destination, une tolérance aux pannes naturelle grâce
à la redondance, et une scalabilité linéaire en $k^3/4$
hôtes.

\section{Sécurité dans les Réseaux SDN}

\subsection{Menaces Spécifiques au SDN}
Le paradigme SDN introduit de nouvelles vulnérabilités
spécifiques \cite{sezer2013ready}. La centralisation du
contrôleur crée un point unique de défaillance qui
constitue une cible privilégiée des attaquants. La
saturation du canal de contrôle par un flood de messages
PacketIn peut paralyser le contrôleur. La manipulation
des règles de flux par injection de règles malveillantes
et les attaques DDoS saturant le réseau représentent
les principales menaces identifiées dans les SDN.

\subsection{Attaques DDoS : Types et Mécanismes}

\begin{table}[H]
\centering
\caption{Types d'attaques DDoS et leurs caractéristiques}
\label{tab:ddos_types}
\begin{tabular}{lll}
\toprule
\textbf{Type} & \textbf{Protocole} & \textbf{Mécanisme} \\
\midrule
SYN flood & TCP & Connexions semi-ouvertes \\
UDP flood & UDP & Saturation bande passante \\
ICMP flood & ICMP & Ping massif \\
HTTP flood & TCP/HTTP & Requêtes web massives \\
Amplification & UDP/DNS & Réflexion et amplification \\
\bottomrule
\end{tabular}
\end{table}

\subsection{Méthodes de Détection DDoS}

\subsubsection{Détection par Entropie de Shannon}
L'entropie de Shannon mesure la diversité d'une
distribution de probabilités. Appliquée aux adresses IP
du trafic réseau, une baisse soudaine de l'entropie
indique une concentration anormale caractéristique d'une
attaque DDoS \cite{braga2010lightweight}. La formule de
l'entropie normalisée est :

\begin{equation}
H_{norm} = \frac{-\sum_{i=1}^{n} p_i \log_2(p_i)}{\log_2(n)}
\label{eq:entropie_norm}
\end{equation}

où $p_i$ désigne la proportion de trafic associée à
l'adresse IP $i$, et $n$ représente le nombre total
d'adresses IP distinctes observées dans l'intervalle
d'analyse.

\subsubsection{Détection par Machine Learning}
Les travaux récents démontrent que les algorithmes
d'apprentissage automatique améliorent considérablement
la précision de détection des attaques DDoS dans les SDN.
Tang et al. \cite{tang2020ddos} obtiennent des taux de
détection supérieurs à 99\% avec l'algorithme Random
Forest appliqué aux statistiques de flux OpenFlow. Xiao
et al. \cite{xiao2021deep} montrent l'efficacité des
réseaux de neurones profonds (LSTM, CNN) pour la détection
en temps réel, tandis qu'Ahuja et al.
\cite{ahuja2021comprehensive} proposent une revue
complète des techniques de détection DDoS dans les SDN,
soulignant les avantages comparatifs des approches
hybrides entropie/ML.

\subsubsection{Détection par Analyse des Flux}
Le nombre de destinations uniques contactées par une
même IP source constitue un indicateur efficace pour
détecter les scans de ports et les attaques distribuées
\cite{mousavi2015early}. Un seuil de 10 destinations
uniques par IP source dans un intervalle de 5 secondes
déclenche une alerte de type scan dans notre implémentation.

\subsection{Mécanismes de Mitigation}
Les travaux de Shin et al. \cite{shin2013avant} sur
AVANT-GUARD montrent qu'une mitigation efficace dans les
SDN doit combiner la rapidité de réaction et la précision
du ciblage pour éviter les blocages injustifiés. Boite
et al. \cite{boite2021toward} proposent une approche
graduée basée sur l'analyse d'anomalies de flux, proche
de notre implémentation à trois niveaux progressifs :
le blocage par règle DROP, la redirection vers un honeypot
pour capturer le comportement de l'attaquant, et le rate
limiting par règles à timeout court.

\subsection{Pare-feu Stateful SDN}
Contrairement aux pare-feux traditionnels qui opèrent
de manière distribuée sur chaque équipement, un pare-feu
stateful SDN centralise le suivi des connexions TCP au
niveau du contrôleur et installe dynamiquement des règles
de filtrage adaptées sur les switches. Les états TCP
gérés dans ce projet sont : SYN (nouvelle connexion
initiée), ESTABLISHED (connexion confirmée par SYN-ACK)
et CLOSING (connexion en cours de fermeture par FIN).

\subsection{Travaux Existants et Positionnement}

\begin{table}[H]
\centering
\caption{Comparaison avec les travaux existants}
\label{tab:related_work}
\begin{tabular}{llll}
\toprule
\textbf{Travail} & \textbf{Méthode} &
\textbf{Mitigation} & \textbf{QoS} \\
\midrule
Braga et al. (2010) \cite{braga2010lightweight} &
Clustering & DROP & Non \\
Shin et al. (2013) \cite{shin2013avant} &
Statistiques & DROP & Non \\
Mousavi et al. (2015) \cite{mousavi2015early} &
Entropie & DROP & Non \\
Tang et al. (2020) \cite{tang2020ddos} &
ML/Random Forest & DROP & Non \\
Boite et al. (2021) \cite{boite2021toward} &
Anomalie flux & Graduée & Non \\
\textbf{Ce projet} & Entropie + seuil &
DROP + Honeypot + RL & Oui \\
\bottomrule
\end{tabular}
\end{table}

Notre approche se distingue par la combinaison de la
détection par entropie avec une mitigation graduée en
trois niveaux et l'intégration simultanée d'un module
QoS, une combinaison qui n'est proposée par aucun des
travaux existants recensés.

\section{Qualité de Service (QoS) dans les SDN}

\subsection{Mécanismes QoS Traditionnels}
Dans les réseaux traditionnels, la QoS est implémentée
via plusieurs mécanismes complémentaires. DiffServ
(Differentiated Services) repose sur le marquage DSCP
des paquets pour les classifier en différentes classes
de service. IntServ (Integrated Services) réserve des
ressources bout-en-bout via le protocole RSVP (Resource
Reservation Protocol). Les files d'attente prioritaires
telles que WFQ (Weighted Fair Queuing), PQ (Priority
Queuing) et CBWFQ (Class-Based WFQ) permettent de
différencier le traitement des flux au niveau de chaque
interface réseau. Ces mécanismes sont configurés
statiquement sur chaque équipement, sans adaptabilité
dynamique aux changements de conditions réseau.

\subsection{QoS Dynamique avec SDN}
Le SDN apporte une révolution dans la gestion de la QoS
en permettant une approche centralisée et dynamique
\cite{liu2022qos}. Le contrôleur classifie les flux et
installe les règles de priorité sur tous les switches
simultanément, sans intervention manuelle. Les politiques
QoS peuvent être modifiées en temps réel en réponse aux
variations de trafic ou aux priorités applicatives. La
priorisation s'effectue au niveau du flux individuel
(5-tuple : IP src/dst, port src/dst, protocole), offrant
une granularité bien supérieure aux mécanismes DiffServ.
Dans ce projet, trois niveaux de QoS sont définis et
appliqués via des règles OpenFlow de priorités 10, 15,
20, 25 et 30.

\subsection{SDN et NFV : Convergence}
Le SDN s'intègre naturellement avec la virtualisation
des fonctions réseau (NFV -- Network Functions
Virtualization). Cette convergence SDN/NFV permet de
déployer des fonctions réseau comme les pare-feux, les
load balancers et les systèmes de détection d'intrusion
sous forme de machines virtuelles ou de conteneurs,
orchestrées dynamiquement par le contrôleur SDN selon
les besoins du trafic.

\section{Vers une Programmation Avancée du Plan de Données}
Le protocole P4 (Programming Protocol-Independent Packet
Processors) \cite{bosshart2014p4} représente une évolution
majeure du paradigme SDN en permettant de programmer
directement le comportement du plan de données des
switches, indépendamment du protocole utilisé. Contrairement
à OpenFlow dont le comportement des switches est fixé par
la spécification du protocole, P4 permet de définir
entièrement les champs analysés, les tables de flux et
les actions appliquées. Cette flexibilité ouvre la voie
à des implémentations de détection DDoS directement
dans le datapath réseau, avec des temps de réaction de
l'ordre de la nanoseconde, sans solliciter le contrôleur.

\section{Conclusion}
Cet état de l'art confirme que le SDN offre des avantages
significatifs sur les réseaux traditionnels en termes de
flexibilité, programmabilité et gestion proactive de la
sécurité. Le choix du contrôleur Ryu, du protocole
OpenFlow 1.3 et de l'émulateur Mininet est pleinement
justifié par leur maturité, leur complémentarité et leur
large adoption dans la communauté académique
\cite{benzekki2016software}. La littérature récente sur
la détection DDoS par machine learning (2020-2025) ouvre
des perspectives d'amélioration directes pour ce projet.
Notre contribution originale réside dans la combinaison
inédite de la détection par entropie avec une mitigation
graduée en trois niveaux et un module QoS intégré.
\newpage

\chapter{Architecture Proposée}

\section{Vue Générale}
L'architecture SDN proposée s'organise en trois couches
conformément au modèle SDN standard, avec quatre modules
applicatifs déployés sur le contrôleur Ryu. Chaque module
est implémenté sous forme d'une application Ryu autonome
pouvant être lancée indépendamment ou en combinaison avec
les autres :
\begin{figure}[H]
\centering
\includegraphics[width=0.9\textwidth]{architecture_sdn.png}
\caption{Architecture globale du système SDN proposé}
\label{fig:architecture}
\end{figure}

\begin{table}[H]
\centering
\caption{Architecture en couches du système développé}
\label{tab:arch}
\begin{tabular}{lll}
\toprule
\textbf{Couche} & \textbf{Composant} & \textbf{Rôle} \\
\midrule
Application & Module QoS, Module Sécurité & Logique métier \\
Contrôle & Contrôleur Ryu (Python) & Gestion centralisée \\
Infrastructure & Switches OVS (Mininet) & Transfert paquets \\
\bottomrule
\end{tabular}
\end{table}

\section{Topologie Réseau}

\subsection{Topologie Arbre (Référence)}
La topologie principale utilisée pour les tests est une
topologie arbre (depth=2, fanout=4). Cette topologie
comprend 1 switch racine (s1), 4 switches intermédiaires
(s2, s3, s4, s5) et 16 hôtes (h1 à h16) dont les
adresses IP s'étendent de 10.0.0.1 à 10.0.0.16. Elle
offre un bon compromis entre complexité de gestion et
représentativité pour les tests de performance et de
sécurité.

\subsection{Topologie Fat-Tree (Avancée)}
Une topologie Fat-Tree k=4 a également été implémentée
dans \texttt{fat\_tree.py}. Contrairement à la topologie
arbre, le Fat-Tree offre des chemins multiples entre
toute paire source-destination, ce qui permet un
reroutage automatique en cas de panne de lien et une
meilleure utilisation de la bande passante disponible.
\begin{figure}[H]
\centering
\includegraphics[width=0.95\textwidth]{topologie_arbre.png}
\caption{Topologie arbre depth=2, fanout=4 -- 
5 switches et 16 hôtes}
\label{fig:topo_arbre}
\end{figure}

\begin{figure}[H]
\centering
\includegraphics[width=0.95\textwidth]{diagramme_modules.png}
\caption{Diagramme des 5 modules et leurs interactions 
avec le contrôleur Ryu}
\label{fig:modules}
\end{figure}

\section{Modules Développés}

\subsection{Module de Surveillance (monitoring.py)}
Ce module assure la collecte périodique des statistiques
OpenFlow de tous les switches connectés toutes les 5
secondes. Il recueille les statistiques par règle de flux
(flow\_stats : paquets, octets, priorité) et les
statistiques par port physique (port\_stats : rx/tx
paquets et octets). Ces métriques alimentent en temps
réel les algorithmes de routage dynamique et de
détection DDoS.

\subsection{Module de Routage Dynamique (routing.py)}
Ce module implémente un routage adaptatif en sélectionnant
le port de sortie optimal selon un score de qualité
calculé pour chaque port disponible :

\begin{equation}
score = taux\_perte(\%) \times 2 + d\acute{e}bit(MB/s)
\label{eq:score}
\end{equation}

Le port affichant le score le plus bas est sélectionné,
car il combine un faible taux de perte et une charge
modérée. Le taux de perte est pondéré deux fois plus
que le débit afin de pénaliser davantage les liens
instables que les liens simplement chargés.

\subsection{Module de Répartition de Charge (load\_balancer.py)}
Ce module distribue les flux réseau équitablement entre
les ports disponibles via l'algorithme Round-Robin.
Chaque nouveau flux (identifié par le couple src\_mac,
dst\_mac) se voit assigner un port de sortie de manière
cyclique :

\begin{equation}
port = ports[compteur \mod nb\_ports]
\label{eq:rr}
\end{equation}

L'assignation est mémorisée pour garantir que tous les
paquets d'un même flux empruntent le même chemin, assurant
ainsi la cohérence et l'ordre de livraison des paquets.

\subsection{Module QoS (qos.py)}
Ce module installe des règles de priorisation sur les
switches dès leur connexion au contrôleur, garantissant
une différenciation de service dès le premier paquet.
Trois niveaux de priorité sont définis :

\begin{table}[H]
\centering
\caption{Politiques QoS implémentées}
\label{tab:qos}
\begin{tabular}{llll}
\toprule
\textbf{Niveau} & \textbf{Priorité OF} &
\textbf{Trafic} & \textbf{Port} \\
\midrule
HAUTE & 30 & iperf3 TCP & 5201 \\
HAUTE & 25 & SSH & 22 \\
MOYENNE & 20 & TCP général & * \\
MOYENNE & 15 & ICMP & * \\
BASSE & 10 & UDP & * \\
\bottomrule
\end{tabular}
\end{table}

\subsection{Module de Sécurité (security\_manager.py)}
Ce module unifié intègre trois sous-systèmes de sécurité
complémentaires qui opèrent de manière coordonnée.

\subsubsection{Détection DDoS}
La détection repose sur trois algorithmes complémentaires
appliqués toutes les 5 secondes : le calcul de l'entropie
de Shannon des distributions d'adresses IP source et
destination, la mesure du taux de paquets par seconde par
IP source, et le comptage des destinations uniques par IP
source pour identifier les scans de ports.

\subsubsection{Mitigation Automatique Graduée}
La mitigation est appliquée progressivement selon la
sévérité de la menace détectée, en trois niveaux :
Rate Limiting (priorité 80, 30s) pour limiter le débit,
Honeypot (priorité 90, 120s) pour rediriger le trafic
vers h16 (10.0.0.16) via SetField, et DROP (priorité
100, 60s) pour bloquer totalement l'IP attaquante.

\subsubsection{Pare-feu Stateful}
Le pare-feu suit les états des connexions TCP (SYN,
ESTABLISHED, CLOSING) et applique des règles ACL
configurables bloquant les protocoles non sécurisés
(Telnet, FTP, RPC, TFTP) tout en autorisant les services
légitimes (SSH, HTTP, HTTPS, iperf3).

\section{Paramètres de Détection DDoS}
Le seuil de \textbf{2 paquets par seconde} a été
délibérément choisi en adéquation avec l'environnement
Mininet simulé. Dans cet environnement, le trafic de
fond généré spontanément par les 16 hôtes de la topologie
reste inférieur à 1 pkt/s par IP source, ce qui rend le
seuil de 2 pkt/s pertinent pour détecter les floods dès
quelques dizaines de paquets par seconde, tout en évitant
les faux positifs. Comme le recommandent Braga et al.
\cite{braga2010lightweight}, ce seuil devrait être
recalibré empiriquement sur au moins 24 heures de trafic
normal avant tout déploiement en environnement réel.

\begin{table}[H]
\centering
\caption{Paramètres de détection DDoS}
\label{tab:params}
\begin{tabular}{lll}
\toprule
\textbf{Paramètre} & \textbf{Valeur} & \textbf{Description} \\
\midrule
SEUIL\_PAQUETS\_PAR\_SEC & 2 & Alerte si $>$ 2 pkt/s \\
SEUIL\_ENTROPIE & 0.5 & Alerte si entropie $<$ 0.5 \\
SEUIL\_FLUX\_PAR\_IP & 10 & Alerte si $>$ 10 destinations \\
INTERVALLE\_ANALYSE & 5s & Période d'analyse \\
\bottomrule
\end{tabular}
\end{table}

\section{Conclusion du Chapitre}
Ce chapitre a présenté l'architecture SDN proposée,
articulée autour de quatre modules complémentaires
déployés sur le contrôleur Ryu : monitoring, routage
dynamique, QoS et sécurité. Cette architecture modulaire
offre une gestion centralisée et programmable du réseau,
avec des mécanismes de sécurité proactifs adaptés aux
menaces contemporaines. Les algorithmes de score de port
et d'entropie de Shannon constituent les éléments
techniques centraux de cette architecture. Le chapitre
suivant détaille l'implémentation de chacun de ces
modules en Python avec Ryu et Mininet.
\newpage

\chapter{Implémentation}

\section{Environnement de Développement}

\subsection{Outils et Technologies}

\begin{table}[H]
\centering
\caption{Environnement de développement utilisé}
\label{tab:outils}
\begin{tabular}{ll}
\toprule
\textbf{Composant} & \textbf{Version/Détail} \\
\midrule
Système d'exploitation & Ubuntu 22.04 LTS \\
Émulateur réseau & Mininet 2.3+ \\
Contrôleur SDN & Ryu 4.34 \\
Langage & Python 3.12 \\
Protocole & OpenFlow 1.3 \\
Switch virtuel & Open vSwitch (OVS) \\
Génération trafic & iperf3, hping3 \\
Monitoring & tshark, Wireshark \\
Versioning & Git / GitHub \\
\bottomrule
\end{tabular}
\end{table}

\subsection{Installation et Configuration}
L'environnement SDN a été déployé dans un environnement
virtuel Python afin d'isoler les dépendances et de
garantir la reproductibilité des expériences sur
différentes machines :

\begin{lstlisting}[language=bash,
caption=Installation de l'environnement SDN]
# Creation de l'environnement virtuel Python
python3 -m venv ~/projet-sdn/venv
source ~/projet-sdn/venv/bin/activate

# Installation du controleur Ryu et dependances
pip install ryu matplotlib numpy

# Lancement du controleur SDN
ryu-manager traffic_manager.py

# Lancement de Mininet avec topologie arbre
sudo mn --controller=remote,ip=127.0.0.1,port=6633 \
        --switch=ovsk,protocols=OpenFlow13 \
        --topo=tree,depth=2,fanout=4 --mac

# Verification de la connectivite initiale
mininet> pingall
\end{lstlisting}

\section{Implémentation du Module de Surveillance}
Le module de monitoring collecte périodiquement les
statistiques OpenFlow via des requêtes OFPPortStatsRequest
toutes les 5 secondes. Pour chaque port de chaque switch,
il calcule le débit instantané en bytes par seconde par
différence entre deux relevés successifs de compteurs,
et le taux de perte de paquets comme rapport entre les
paquets perdus et les paquets émis :

\begin{lstlisting}[language=Python,
caption=Calcul du débit et taux de perte (traffic\_manager.py)]
def _port_stats_reply_handler(self, ev):
    dpid = ev.msg.datapath.id
    for stat in ev.msg.body:
        port = stat.port_no
        # Calcul du debit instantane (bytes/sec)
        curr = stat.tx_bytes + stat.rx_bytes
        speed = (curr - self.prev_bytes[dpid][port]) / 5.0
        self.prev_bytes[dpid][port] = curr
        # Calcul du taux de perte de paquets
        loss = 0.0
        if stat.tx_packets > 0:
            loss = max(0.0,
                (stat.tx_packets - stat.rx_packets)
                / stat.tx_packets * 100)
        self.port_stats[dpid][port] = {
            'speed': speed, 'loss': loss}
\end{lstlisting}

\section{Implémentation du Routage Dynamique}
Le module de routage sélectionne le port de sortie
optimal en calculant un score pour chaque port disponible.
Ce score combine le taux de perte (pondéré ×2) et le
débit normalisé en MB/s, favorisant les liens à la fois
fiables et peu chargés :

\begin{lstlisting}[language=Python,
caption=Sélection du meilleur port (routing.py)]
def _best_port(self, dpid, available_ports):
    best_port = None
    best_score = float('inf')
    for port in available_ports:
        metrics = self.port_stats[dpid].get(port, {})
        speed = metrics.get('speed', 0)
        loss = metrics.get('loss', 0)
        # Score : perte ponderee x2 + charge normalisee
        score = loss * 2 + (speed / 1e6)
        if score < best_score:
            best_score = score
            best_port = port
    return best_port
\end{lstlisting}

\section{Implémentation du Load Balancer}
L'algorithme Round-Robin assigne cycliquement un port
de sortie à chaque nouveau flux, en mémorisant
l'assignation pour garantir la cohérence des paquets :

\begin{lstlisting}[language=Python,
caption=Algorithme Round-Robin (load\_balancer.py)]
def _assign_port_rr(self, dpid, src, dst, available_ports):
    key = (src, dst)
    if key not in self.flow_port[dpid]:
        # Selection cyclique du port
        idx = self.rr_counter[dpid] % len(available_ports)
        self.rr_counter[dpid] += 1
        self.flow_port[dpid][key] = available_ports[idx]
        self.port_load[dpid][available_ports[idx]] += 1
    return self.flow_port[dpid][key]
\end{lstlisting}

\section{Implémentation du Module QoS}
Les règles QoS sont installées sur les switches via
FlowMod dès leur connexion au contrôleur, garantissant
une priorisation du trafic dès le premier paquet reçu :

\begin{lstlisting}[language=Python,
caption=Installation des règles QoS (qos.py)]
def _install_qos_rules(self, datapath):
    parser = datapath.ofproto_parser
    ofproto = datapath.ofproto
    actions = [parser.OFPActionOutput(ofproto.OFPP_NORMAL)]

    # HAUTE priorite (30) : iperf3 TCP port 5201
    match = parser.OFPMatch(eth_type=0x0800,
                            ip_proto=6, tcp_dst=5201)
    self.add_flow(datapath, 30, match, actions)

    # HAUTE priorite (25) : SSH TCP port 22
    match = parser.OFPMatch(eth_type=0x0800,
                            ip_proto=6, tcp_dst=22)
    self.add_flow(datapath, 25, match, actions)

    # MOYENNE priorite (15) : ICMP
    match = parser.OFPMatch(eth_type=0x0800, ip_proto=1)
    self.add_flow(datapath, 15, match, actions)

    # BASSE priorite (10) : UDP
    match = parser.OFPMatch(eth_type=0x0800, ip_proto=17)
    self.add_flow(datapath, 10, match, actions)
\end{lstlisting}

\section{Implémentation du Module de Sécurité}

\subsection{Détection DDoS par Entropie de Shannon}
L'algorithme calcule l'entropie de Shannon normalisée
des distributions d'adresses IP source et destination
pour chaque intervalle d'analyse de 5 secondes. Une
valeur proche de 0 déclenche une alerte DDoS :

\begin{lstlisting}[language=Python,
caption=Calcul de l'entropie de Shannon (ddos\_detector.py)]
def _calcul_entropie(self, compteurs):
    total = sum(compteurs.values())
    if total == 0:
        return 1.0  # Pas de trafic = entropie maximale
    entropie = 0.0
    for count in compteurs.values():
        p = count / total
        if p > 0:
            entropie -= p * math.log2(p)
    # Normalisation par log2(N) pour obtenir [0,1]
    nb = len(compteurs)
    if nb > 1:
        entropie /= math.log2(nb)
    return entropie
\end{lstlisting}

\subsection{Mitigation Automatique par Règle DROP}
Lorsqu'une IP source est identifiée comme attaquante,
une règle OpenFlow de haute priorité est installée avec
une liste d'actions vide (DROP) et un hard\_timeout de
60 secondes, assurant un blocage temporaire :

\begin{lstlisting}[language=Python,
caption=Blocage DROP d'une IP suspecte (ddos\_detector.py)]
def _block_ip(self, datapath, ip_src):
    # Verification : IP deja bloquee sur ce switch ?
    if ip_src in self.blocked_ips[datapath.id]:
        return
    parser = datapath.ofproto_parser
    # Correspondance sur l'adresse IP source IPv4
    match = parser.OFPMatch(eth_type=0x0800,
                            ipv4_src=ip_src)
    # Action DROP : liste vide, priorite 100, timeout 60s
    self.add_flow(datapath, 100, match, [], hard=60)
    self.blocked_ips[datapath.id].add(ip_src)
\end{lstlisting}

\subsection{Pare-feu Stateful}
Le pare-feu suit dynamiquement les états des connexions TCP :

\begin{table}[H]
\centering
\caption{États TCP suivis par le pare-feu stateful}
\label{tab:tcp_states}
\begin{tabular}{lll}
\toprule
\textbf{Drapeau TCP} & \textbf{État} & \textbf{Action} \\
\midrule
SYN & Nouvelle connexion & Vérification ACL \\
SYN-ACK & Connexion établie & ESTABLISHED \\
FIN & Fermeture & CLOSING \\
Données & Connexion active & Mise à jour timestamp \\
\bottomrule
\end{tabular}
\end{table}

\subsection{Règles ACL du Pare-feu}

\begin{table}[H]
\centering
\caption{Règles ACL du pare-feu stateful}
\label{tab:acl}
\begin{tabular}{llll}
\toprule
\textbf{Action} & \textbf{Protocole} &
\textbf{Port} & \textbf{Service} \\
\midrule
DENY & TCP & 23 & Telnet \\
DENY & TCP & 21 & FTP \\
DENY & TCP & 135 & RPC \\
DENY & UDP & 69 & TFTP \\
ALLOW & TCP & 22 & SSH \\
ALLOW & TCP & 80 & HTTP \\
ALLOW & TCP & 443 & HTTPS \\
ALLOW & TCP & 5201 & iperf3 \\
ALLOW & ICMP & * & Ping \\
ALLOW & * & * & Défaut \\
\bottomrule
\end{tabular}
\end{table}

\section{Captures Wireshark des Messages OpenFlow}
La surveillance des échanges OpenFlow entre le contrôleur
Ryu et les switches a été réalisée avec Wireshark sur
l'interface loopback avec le filtre \texttt{openflow\_v4}.
Les messages observés comprennent les requêtes de
statistiques OFPT\_MULTIPART\_REQUEST envoyées par Ryu
toutes les 5 secondes, les réponses OFPT\_MULTIPART\_REPLY
des switches avec les compteurs de ports et flux, les
messages keepalive OFPT\_ECHO\_REPLY pour maintenir la
connexion, et les messages OFPT\_FLOW\_MOD pour
l'installation des règles de flux. Au total, 30 messages
OpenFlow ont été échangés en 10 secondes, confirmant un
overhead de contrôle minimal de 3 messages par seconde.

\section{Conclusion du Chapitre}
Ce chapitre a présenté l'implémentation complète des
quatre modules de l'architecture SDN en Python avec Ryu.
Les extraits de code illustrent les algorithmes techniques
clés : le calcul de l'entropie de Shannon pour la
détection DDoS, la sélection du meilleur port par score
pour le routage adaptatif, l'algorithme Round-Robin pour
le load balancing, et l'installation des règles QoS et
de sécurité via FlowMod OpenFlow. La section suivante
évalue les performances de cette implémentation sur
quatre scénarios de test distincts.
\newpage

\chapter{Tests et Évaluation des Performances}

\section{Environnement de Test}
Les tests ont été réalisés sur la topologie arbre
(depth=2, fanout=4) avec 16 hôtes et 5 switches OVS
connectés au contrôleur Ryu via OpenFlow 1.3. Les outils
de mesure utilisés sont ping (latence et taux de perte),
iperf3 (débit TCP/UDP et jitter), hping3 (génération
d'attaques DDoS) et tshark/Wireshark (capture et analyse
des messages OpenFlow).

\section{Méthodologie de Comparaison SDN vs Réseau Traditionnel}
Afin de garantir une comparaison scientifiquement valide,
le réseau traditionnel a été simulé dans \textbf{le même
environnement Mininet} avec la même topologie arbre
(depth=2, fanout=4) et les mêmes paramètres de bande
passante. La seule différence introduite est l'absence
de contrôleur SDN, obtenue par l'option
\texttt{--controller=none}. Dans cette configuration,
les switches OVS fonctionnent en mode autonome avec
un simple apprentissage MAC basique, sans optimisation
de routage, sans politiques QoS et sans aucun mécanisme
de sécurité. Les mêmes outils de mesure (iperf3, ping,
hping3) ont été utilisés dans les deux configurations
avec des paramètres identiques, afin de garantir la
comparabilité des résultats obtenus.

\section{Scénario 1 : Trafic Normal}

\subsection{Description}
Ce scénario évalue les performances du réseau SDN en
conditions normales, avec trois niveaux de charge
croissants : faible (1 flux TCP), moyenne (4 flux
simultanés) et forte (8 flux simultanés).

\subsection{Résultats}

\begin{table}[H]
\centering
\caption{Résultats Scénario 1 -- Trafic normal}
\label{tab:scenario1}
\begin{tabular}{lllll}
\toprule
\textbf{Charge} & \textbf{Flux} & \textbf{Débit TCP} &
\textbf{Latence avg} & \textbf{Perte} \\
\midrule
Faible & 1 flux & 23,5 Gbps & 0,134 ms & 0\% \\
Moyenne & 4 flux & 24,2 Gbps & -- & 0\% \\
Forte & 8 flux & 21,5 Gbps & -- & 0\% \\
\bottomrule
\end{tabular}
\end{table}

\begin{table}[H]
\centering
\caption{Métriques UDP et Jitter -- Scénario 1}
\label{tab:jitter}
\begin{tabular}{llll}
\toprule
\textbf{Débit UDP} & \textbf{Jitter} &
\textbf{Perte UDP} & \textbf{Paquets} \\
\midrule
987 Mbits/sec & 0,001 ms & 1,2\% & 863 247 \\
\bottomrule
\end{tabular}
\end{table}

\begin{figure}[H]
\centering
\includegraphics[width=0.9\textwidth]{graphique1_debit.png}
\caption{Débit TCP selon la charge -- SDN vs Réseau Traditionnel}
\label{fig:debit}
\end{figure}

\subsection{Analyse}
Les résultats démontrent d'excellentes performances du
réseau SDN en conditions normales. Le débit TCP se
maintient entre 21,5 et 24,2 Gbps quelle que soit la
charge appliquée, ce qui confirme la capacité du routage
dynamique à préserver des performances élevées sous charge
variable. Le jitter UDP de 0,001 ms témoigne d'une très
faible variabilité de la latence, une qualité essentielle
pour les applications temps réel et la téléphonie sur IP.

\section{Scénario 2 : Attaque DDoS}

\subsection{Description}
Ce scénario évalue l'impact des attaques DDoS sur les
performances réseau et l'efficacité du système de
détection et mitigation automatique. Deux types d'attaques
ont été générés depuis h1 : un ICMP flood par envoi
rapide de pings (\texttt{ping -i 0.01}) et un SYN+UDP
flood combiné via hping3 (\texttt{hping3 -S -p 80
--flood} et \texttt{hping3 -2 -p 53 --flood}).

\subsection{Résultats}

\begin{table}[H]
\centering
\caption{Impact des attaques DDoS sur les performances}
\label{tab:scenario2}
\begin{tabular}{llll}
\toprule
\textbf{Phase} & \textbf{Latence avg} &
\textbf{Latence max} & \textbf{Perte} \\
\midrule
Avant attaque & 0,478 ms & 1,966 ms & 0\% \\
ICMP flood & 4,793 ms & 45,378 ms & 0\% \\
SYN+UDP flood & 12 718 ms & 16 833 ms & 10\% \\
Après mitigation & 2,389 ms & 28,498 ms & 0\% \\
\bottomrule
\end{tabular}
\end{table}

\begin{table}[H]
\centering
\caption{Performances du système de détection/mitigation}
\label{tab:detection}
\begin{tabular}{ll}
\toprule
\textbf{Métrique} & \textbf{Valeur} \\
\midrule
Temps de détection & 5 secondes \\
Temps de mitigation & $<$ 1 seconde \\
Débit attaque détecté & 446 pkt/s \\
Seuil de déclenchement & 2 pkt/s \\
Durée du blocage DROP & 60 secondes \\
\bottomrule
\end{tabular}
\end{table}

\begin{figure}[H]
\centering
\includegraphics[width=0.9\textwidth]{graphique2_ddos.png}
\caption{Impact des attaques DDoS sur la latence et le taux de perte}
\label{fig:ddos}
\end{figure}

\subsection{Analyse}
Le système de détection identifie les attaques en un
seul intervalle d'analyse de 5 secondes et déclenche la
mitigation automatique en moins d'une seconde après
détection. Après l'installation de la règle DROP sur
l'IP attaquante, la latence revient à 2,389 ms, contre
0,478 ms avant l'attaque. Cette légère différence résiduelle
s'explique par la congestion accumulée pendant les
5 secondes précédant la détection.

\section{Scénario 3 : Panne d'un Lien}

\subsection{Description}
Ce scénario simule la panne physique du lien entre les
switches s1 et s2 par la commande Mininet
\texttt{link s1 s2 down}, puis son rétablissement par
\texttt{link s1 s2 up}. Ce test évalue la résilience
du réseau SDN face aux défaillances d'infrastructure.

\subsection{Résultats}

\begin{table}[H]
\centering
\caption{Impact de la panne du lien s1-s2}
\label{tab:scenario3}
\begin{tabular}{lll}
\toprule
\textbf{Phase} & \textbf{Latence avg} & \textbf{Perte} \\
\midrule
Avant panne & 5,056 ms & 0\% \\
Pendant panne & -- & 100\% \\
Après rétablissement & 6,086 ms & 0\% \\
\bottomrule
\end{tabular}
\end{table}

\begin{figure}[H]
\centering
\includegraphics[width=0.9\textwidth]{graphique3_panne.png}
\caption{Impact de la panne du lien s1-s2 sur la connectivité}
\label{fig:panne}
\end{figure}

\subsection{Analyse}
La topologie arbre ne disposant que d'un seul chemin
entre h1 et h16 via le lien s1-s2, la panne entraîne
une interruption totale de la connectivité (100\% de
perte). Après rétablissement du lien, la connectivité
est restaurée avec une latence légèrement augmentée
(6,086 ms vs 5,056 ms avant panne), due au délai
nécessaire à la réinstallation des règles de flux par
le contrôleur. La topologie Fat-Tree k=4 permettrait
quant à elle un reroutage automatique vers les chemins
alternatifs disponibles, sans interruption de service.

\section{Scénario 4 : QoS sous Congestion}

\subsection{Description}
Ce scénario évalue l'application des politiques QoS en
présence d'un trafic mixte simultané : un flux TCP iperf3
en haute priorité et un flux UDP 100M en basse priorité,
afin de vérifier que le trafic critique est bien
préservé en situation de congestion.

\subsection{Résultats}

\begin{table}[H]
\centering
\caption{Résultats des politiques QoS}
\label{tab:scenario4}
\begin{tabular}{llll}
\toprule
\textbf{Trafic} & \textbf{Priorité} &
\textbf{Débit} & \textbf{Latence} \\
\midrule
iperf3 TCP 5201 & HAUTE (30) & 19,7 Gbps & -- \\
SSH TCP 22 & HAUTE (25) & Variable & Faible \\
TCP général & MOYENNE (20) & Variable & Normale \\
ICMP ping & MOYENNE (15) & -- & 0,219 ms \\
UDP 100M & BASSE (10) & 100 Mbps & -- \\
\bottomrule
\end{tabular}
\end{table}

\subsection{Analyse}
Les politiques QoS fonctionnent conformément aux
attentes. Le trafic TCP iperf3 obtient 19,7 Gbps en
haute priorité, tandis que l'ICMP maintient une latence
de 0,219 ms malgré la congestion générée par le flux
UDP concurrent. Le trafic UDP est bien relégué en basse
priorité, démontrant l'efficacité de la différenciation
de service par règles OpenFlow directement sur les switches.

\section{Overhead du Contrôleur}

\begin{table}[H]
\centering
\caption{Overhead du contrôleur Ryu}
\label{tab:overhead}
\begin{tabular}{ll}
\toprule
\textbf{Métrique} & \textbf{Valeur} \\
\midrule
CPU (sous charge DDoS) & 21,5\% \\
Mémoire utilisée & 65 MB (0,8\%) \\
Messages OpenFlow & 30 messages / 10s \\
Temps de détection DDoS & 5 secondes \\
Temps de mitigation & $<$ 1 seconde \\
\bottomrule
\end{tabular}
\end{table}

\begin{figure}[H]
\centering
\includegraphics[width=0.9\textwidth]{graphique5_overhead.png}
\caption{Overhead CPU, mémoire et messages OpenFlow du contrôleur Ryu}
\label{fig:overhead}
\end{figure}

\section{Comparaison SDN vs Réseau Traditionnel}

\begin{table}[H]
\centering
\caption{Comparaison SDN vs Réseau Traditionnel}
\label{tab:comparaison}
\begin{tabular}{llll}
\toprule
\textbf{Critère} & \textbf{SDN (Ryu)} &
\textbf{Traditionnel} & \textbf{Avantage} \\
\midrule
Latence avg (normal) & 0,134 ms & 0,279 ms & SDN \\
Débit TCP (1 flux) & 23,5 Gbps & 23,9 Gbps & Comparable \\
Débit TCP (4 flux) & 24,2 Gbps & 22,0 Gbps & SDN \\
Jitter UDP & 0,001 ms & Non mesuré & SDN \\
Perte paquets & 0\% & 0\% & Égal \\
Détection DDoS & 5 secondes & Aucune & SDN \\
Mitigation DDoS & Auto ($<$1s) & Manuelle & SDN \\
QoS dynamique & 3 niveaux & Non & SDN \\
Reroutage après panne & Partiel & Statique & SDN \\
\bottomrule
\end{tabular}
\end{table}
\subsection{Analyse Comparative}
Le réseau SDN présente une latence 52\% inférieure
(0,134 ms vs 0,279 ms) en conditions normales. Le débit
TCP est comparable pour un flux unique mais nettement
supérieur sous forte charge (24,2 Gbps vs 22,0 Gbps
pour 4 flux simultanés), grâce au routage dynamique et
au load balancing. L'avantage décisif du SDN réside dans
la gestion de la sécurité : la détection automatique
des attaques en 5 secondes et la mitigation en moins
d'une seconde sont impossibles à réaliser dans un réseau
traditionnel sans infrastructure dédiée et coûteuse.
\newpage
\section{Comparaison Topologie Arbre vs Fat-Tree k=4}
La topologie Fat-Tree k=4, implémentée dans \texttt{fat\_tree.py},
a été démarrée avec succès sous Mininet avec 20 switches
(4 core, 8 agrégation, 8 edge) et 16 hôtes correctement
créés. Les tests préliminaires montrent cependant une
connectivité partielle due au temps de convergence du
protocole STP (Spanning Tree Protocol), inhérent aux
topologies à chemins multiples.

\begin{table}[H]
\centering
\caption{Comparaison Topologie Arbre vs Fat-Tree k=4}
\label{tab:fattree}
\begin{tabular}{lll}
\toprule
\textbf{Critère} & \textbf{Arbre (depth=2)} &
\textbf{Fat-Tree (k=4)} \\
\midrule
Nombre de chemins & 1 (unique) & $(k/2)^2 = 4$ \\
Tolérance aux pannes & Aucune & Reroutage possible \\
Bande passante bisect. & Limitée & Maximale \\
Nombre de switches & 5 & 20 \\
Nombre d'hôtes & 16 & 16 \\
Convergence STP & N/A & 60--90 secondes \\
Connectivité obtenue & 100\% & Partielle \\
\bottomrule
\end{tabular}
\end{table}

La connectivité partielle observée s'explique par la
complexité du protocole STP sur 20 switches interconnectés.
Une implémentation complète nécessiterait un algorithme
de routage multipath dédié, comme ECMP (Equal-Cost
Multi-Path) \cite{alfares2008scalable}, ce qui constitue
une perspective d'amélioration directe de ce projet.
\section{Conclusion du Chapitre}
Ce chapitre a présenté l'évaluation complète de l'architecture
SDN proposée à travers quatre scénarios de test distincts.
Les résultats obtenus confirment la supériorité du SDN sur
le réseau traditionnel en termes de latence (réduction de 52\%),
de sécurité (détection DDoS en 5 secondes, mitigation en moins
d'une seconde) et de gestion de la qualité de service (trois
niveaux de priorité fonctionnels). La méthodologie comparative
rigoureuse, basée sur le même environnement Mininet avec et
sans contrôleur SDN, garantit la validité scientifique des
résultats présentés. Les perspectives d'amélioration identifiées,
notamment l'évaluation sur topologie Fat-Tree et l'intégration
du Machine Learning pour la détection DDoS, feront l'objet
de travaux futurs.
\newpage

\chapter{Conclusion et Perspectives}

\section{Bilan du Projet}
Ce projet a permis de concevoir et de déployer une
architecture SDN complète répondant aux exigences de
gestion dynamique du trafic et de sécurité réseau.
Les six phases du projet ont été menées à bien : étude
bibliographique et état de l'art, déploiement de
l'environnement Mininet et Ryu, développement des modules
de gestion du trafic (monitoring, routage, QoS, load
balancing), développement du module de sécurité (détection
DDoS, mitigation, pare-feu stateful), tests et évaluation
sur quatre scénarios, et rédaction du présent rapport.

\section{Résultats Obtenus}

\begin{table}[H]
\centering
\caption{Récapitulatif des résultats obtenus}
\label{tab:resultats}
\begin{tabular}{lll}
\toprule
\textbf{Module} & \textbf{Métrique} & \textbf{Résultat} \\
\midrule
Trafic normal & Débit TCP & 23,5 Gbps \\
Trafic normal & Latence avg & 0,134 ms \\
Trafic normal & Jitter UDP & 0,001 ms \\
Trafic normal & Perte paquets & 0\% \\
Sécurité & Détection DDoS & 5 secondes \\
Sécurité & Mitigation & $<$ 1 seconde \\
QoS & Débit TCP prioritaire & 19,7 Gbps \\
QoS & Latence ICMP & 0,219 ms \\
Comparaison & Réduction latence & 52\% vs traditionnel \\
Overhead & CPU contrôleur & 21,5\% \\
Overhead & Mémoire & 65 MB \\
\bottomrule
\end{tabular}
\end{table}

\section{Contributions}

\begin{enumerate}
    \item Un \textbf{module de détection DDoS} original
    combinant trois algorithmes complémentaires : le calcul
    de l'entropie de Shannon des distributions IP, la
    mesure du taux de paquets par seconde par IP source
    et la détection de scans par comptage des destinations.
    \item Un \textbf{système de mitigation automatique
    graduée} en trois niveaux progressifs (rate limiting,
    honeypot, DROP) permettant une réponse proportionnelle
    à la sévérité de l'attaque détectée.
    \item Un \textbf{pare-feu stateful SDN} avec suivi
    des états TCP (SYN, ESTABLISHED, CLOSING) et règles
    ACL configurables, entièrement implémenté au niveau
    du contrôleur Ryu.
    \item Un \textbf{module QoS dynamique} à trois niveaux
    de priorité (HAUTE, MOYENNE, BASSE) installé via des
    règles OpenFlow FlowMod directement sur les switches.
    \item Une \textbf{évaluation comparative complète}
    sur quatre scénarios démontrant la supériorité du
    SDN sur le réseau traditionnel en termes de latence,
    sécurité et gestion du trafic.
\end{enumerate}

\section{Difficultés Rencontrées}

\begin{itemize}
    \item \textbf{Compatibilité Python 3.12/Ryu} : Ryu
    étant initialement conçu pour Python 2 et 3.6,
    des avertissements de dépréciation liés à Eventlet
    ont nécessité le recours à un environnement virtuel
    dédié pour isoler les dépendances.
    \item \textbf{Calibrage des seuils DDoS} : Le seuil
    de détection de 2 pkt/s a dû être ajusté
    empiriquement pour éviter les faux positifs dans
    l'environnement simulé Mininet, où le trafic de
    fond est naturellement faible.
    \item \textbf{Règles OpenFlow conflictuelles} : La
    gestion des priorités entre les règles de sécurité
    (priorité 100) et les règles de routage normal
    (priorité 1) a nécessité une conception soigneuse
    pour éviter les conflits et les comportements
    inattendus.
    \item \textbf{Idle timeout} : Le paramètre idle=2s
    dans ddos\_detector.py s'est avéré nécessaire pour
    forcer la remontée régulière des paquets vers le
    contrôleur et maintenir les compteurs DDoS à jour.
\end{itemize}

\section{Perspectives}

\subsection{Court Terme}
À court terme, plusieurs améliorations directes sont
envisageables. Le développement d'une interface web de
supervision en temps réel (Flask/Django) permettrait de
visualiser les métriques réseau, les alertes DDoS et les
règles de flux actives. L'implémentation d'un mécanisme
de reroutage automatique lors d'une panne de lien
exploiterait les chemins alternatifs disponibles sur la
topologie Fat-Tree. L'évaluation des performances sur
la topologie Fat-Tree k=8 fournirait également des
résultats plus représentatifs d'un environnement
datacenter réel.

\subsection{Moyen Terme}
À moyen terme, l'intégration d'algorithmes de Machine
Learning tels que SVM, Random Forest ou LSTM améliorerait
la précision de la détection DDoS tout en réduisant les
faux positifs, conformément aux résultats encourageants
de Tang et al. \cite{tang2020ddos} et Xiao et al.
\cite{xiao2021deep}. Une étude de scalabilité avec 50
à 500 switches permettrait d'évaluer les limites de
passage à l'échelle du contrôleur Ryu. Le déploiement
d'une architecture multi-contrôleurs Ryu renforcerait
quant à lui la haute disponibilité du système.

\subsection{Long Terme}
À long terme, la migration vers le contrôleur ONOS
permettrait un déploiement en environnement de production
avec une haute disponibilité garantie. L'extension de
l'architecture pour la gestion du trafic IoT avec des
politiques de sécurité adaptées aux contraintes des
objets connectés représente une perspective prometteuse
\cite{abubakar2021iot}. Enfin, l'exploration du langage
P4 permettrait d'implémenter la détection DDoS
directement dans le datapath réseau des switches, avec
des temps de réaction de l'ordre de la nanoseconde
\cite{bosshart2014p4}.

\section{Conclusion Générale}
Ce projet a démontré la faisabilité et l'efficacité d'une
architecture SDN pour la gestion dynamique du trafic et
la sécurité réseau dans un environnement académique. Le
contrôleur Ryu, combiné à l'émulateur Mininet, offre un
environnement de développement et de test idéal pour
prototyper et valider des solutions réseau innovantes à
moindre coût, sans nécessiter d'infrastructure physique
dédiée.

Les résultats obtenus confirment les avantages du SDN :
une latence réduite de 52\% par rapport au réseau
traditionnel (0,134 ms vs 0,279 ms), une détection DDoS
automatique en 5 secondes avec mitigation en moins d'une
seconde, et une gestion QoS dynamique garantissant la
qualité de service pour les applications critiques
même en situation de congestion.

Dans le contexte africain et particulièrement au Sénégal,
le déploiement d'architectures SDN dans les universités,
les centres de données et les opérateurs télécoms
représente une opportunité stratégique majeure pour
moderniser les infrastructures numériques. Le programme
``Sénégal Numérique 2025'' place le développement
numérique au cœur des priorités nationales, et le SDN
peut jouer un rôle déterminant dans cet effort de
modernisation. La programmabilité et la centralisation
du contrôle offrent une flexibilité particulièrement
adaptée aux contraintes des réseaux africains, notamment
la gestion optimale d'une bande passante internationale
encore limitée et coûteuse, et la sécurisation contre
des cyberattaques en croissance constante sur le continent
\cite{ndiaye2022sdn}. Ce travail constitue ainsi une
contribution concrète à la réflexion sur la modernisation
des infrastructures réseau en Afrique de l'Ouest.

Le code source complet, documenté et commenté, est
disponible sur GitHub à l'adresse :
\url{https://github.com/thioraminata753/projet-sdn}.

\newpage

\bibliographystyle{plain}
\bibliography{references}

\appendix

\chapter{Guide d'Installation et de Déploiement}

\section{Prérequis}
\begin{itemize}
    \item Ubuntu 22.04 LTS (ou Debian 11+)
    \item Python 3.12+
    \item Mininet 2.3+
    \item Open vSwitch 2.13+
    \item Git, iperf3, hping3, tshark, Wireshark
\end{itemize}

\section{Installation Pas à Pas}

\begin{lstlisting}[language=bash,
caption=Installation complète de l'environnement]
# Etape 1 : Cloner le depot GitHub
git clone https://github.com/thioraminata753/projet-sdn
cd projet-sdn

# Etape 2 : Creer l'environnement virtuel Python
python3 -m venv venv
source venv/bin/activate

# Etape 3 : Installer les dependances
pip install ryu matplotlib numpy

# Etape 4 : Verifier l'installation
ryu-manager --version
mn --version
iperf3 --version
\end{lstlisting}

\section{Lancement des Modules}

\begin{lstlisting}[language=bash,
caption=Lancement du projet SDN]
# Terminal 1 : Activer venv et lancer Ryu
source ~/projet-sdn/venv/bin/activate
ryu-manager security_manager.py

# Terminal 2 : Lancer Mininet
sudo mn --controller=remote,ip=127.0.0.1,port=6633 \
        --switch=ovsk,protocols=OpenFlow13 \
        --topo=tree,depth=2,fanout=4 --mac

# Terminal 3 : Surveiller les logs
tail -f /tmp/sdn_logs/security_manager.log
\end{lstlisting}

\section{Tests de Validation}

\begin{lstlisting}[language=bash,
caption=Commandes de test dans Mininet]
mininet> pingall
mininet> h1 ping -c 100 10.0.0.16
mininet> h1 iperf3 -s &
mininet> h16 iperf3 -c 10.0.0.1 -t 10
mininet> h16 iperf3 -c 10.0.0.1 -u -b 1G -t 10
mininet> h1 hping3 -S -p 80 --flood 10.0.0.16 &
mininet> link s1 s2 down
mininet> link s1 s2 up
python3 graphiques_phase5.py
\end{lstlisting}

\chapter{Structure du Dépôt GitHub}

Le code source complet est disponible à l'adresse :
\url{https://github.com/thioraminata753/projet-sdn}

\begin{table}[H]
\centering
\caption{Structure du dépôt GitHub}
\label{tab:github}
\begin{tabular}{ll}
\toprule
\textbf{Fichier/Dossier} & \textbf{Description} \\
\midrule
traffic\_manager.py & Module principal gestion trafic \\
ddos\_detector.py & Détection DDoS par entropie \\
firewall.py & Pare-feu stateful ACL \\
security\_manager.py & Module sécurité unifié \\
qos.py & Gestion QoS 3 niveaux \\
monitoring.py & Surveillance trafic \\
routing.py & Routage dynamique \\
load\_balancer.py & Load balancing Round-Robin \\
mitigation.py & Mitigation automatique \\
fat\_tree.py & Topologie Fat-Tree k=4 \\
graphiques\_phase5.py & Génération graphiques \\
tests/ & Scripts tests automatisés \\
captures/ & Captures d'écran des tests \\
resultats\_graphiques/ & Graphiques matplotlib \\
rapports/ & Rapport intermédiaire PDF \\
\bottomrule
\end{tabular}
\end{table}

\chapter{Paramètres de Configuration}

\begin{table}[H]
\centering
\caption{Paramètres configurables des modules}
\label{tab:params_config}
\begin{tabular}{lll}
\toprule
\textbf{Paramètre} & \textbf{Valeur} & \textbf{Module} \\
\midrule
SEUIL\_PAQUETS\_PAR\_SEC & 2 & ddos\_detector.py \\
SEUIL\_ENTROPIE & 0.5 & ddos\_detector.py \\
SEUIL\_FLUX\_PAR\_IP & 10 & ddos\_detector.py \\
INTERVALLE\_ANALYSE & 5s & ddos\_detector.py \\
HONEYPOT\_IP & 10.0.0.16 & security\_manager.py \\
DROP\_TIMEOUT & 60s & ddos\_detector.py \\
HONEYPOT\_TIMEOUT & 120s & mitigation.py \\
RATE\_LIMIT\_TIMEOUT & 30s & mitigation.py \\
idle\_timeout (routage) & 30s & traffic\_manager.py \\
idle\_timeout (DDoS) & 2s & ddos\_detector.py \\
Port contrôleur & 6633 & Tous les modules \\
\bottomrule
\end{tabular}
\end{table}

\end{document}